% One-time imported draft. Editable native source; no Markdown regeneration.
\section{Step 22 Compare incremental and rational
implementations}\label{007-bresenham-line-s22-step-22-compare-incremental-and-rational-implementations}

The alternative computes each minor coordinate by integer rounded
division along the major axis. The reference updates a running error.
Both must match the same canonical tie rule and full-frame contract.

\subsection{Write or
inspect}\label{007-bresenham-line-s22-write-or-inspect}

Compare the two methods and their measured default instruction totals.

\subsection{Expected
observation}\label{007-bresenham-line-s22-expected-observation}

The reference uses 145 completed instructions for the default fixture;
the rational alternative uses 138. These are simulated instruction
counts, not host timings or cycle estimates.

\subsection{Check your
understanding}\label{007-bresenham-line-s22-check-your-understanding}

What can those totals establish?

\begin{itemize}
\tightlist
\item
  \textbf{A} How many simulated instructions those exact runs completed.
\item
  \textbf{B} That division costs one real hardware cycle.
\item
  \textbf{C} That the alternative is faster on every input and host.
\end{itemize}
